\documentclass[11pt,a4paper]{article}

% -----------------------------------------------------------------------------
% CS6868 Research Mini-Project Report
% -----------------------------------------------------------------------------

\usepackage[utf8]{inputenc}
\usepackage[T1]{fontenc}
\usepackage[margin=1in]{geometry}
\usepackage{microtype}
\usepackage{lmodern}
\usepackage{amsmath,amssymb}
\usepackage{graphicx}
\usepackage{float}
\usepackage{booktabs}
\usepackage{array}
\usepackage{enumitem}
\usepackage{xcolor}
\usepackage{listings}
\usepackage{hyperref}
\usepackage[capitalise,noabbrev]{cleveref}

\hypersetup{
  colorlinks=true,
  linkcolor=black,
  citecolor=blue!60!black,
  urlcolor=blue!60!black,
}

\definecolor{codegray}{gray}{0.95}
\lstset{
  basicstyle=\ttfamily\small,
  backgroundcolor=\color{codegray},
  frame=single,
  framesep=4pt,
  breaklines=true,
  columns=fullflexible,
  keepspaces=true,
  showstringspaces=false,
}

\title{%
  \textbf{CS6868 Research Mini-Project} \\[0.3em]
  \Large Concurrent Union-Find in OCaml 5 Using CAS Atomics%
}
\author{%
  Yaswanth Srinivas \\ \texttt{cs23b047@smail.iitm.ac.in}
  \and
  Prince \\ \texttt{cs23b045@smail.iitm.ac.in}
  \and
  Mukesh \\ \texttt{cs23b040@smail.iitm.ac.in}
}
\date{\today}

\begin{document}
\maketitle

\begin{abstract}
We implement and evaluate a lock-free-style concurrent union-find (disjoint-set
union, DSU) in OCaml~5 using \texttt{Atomic.compare\_and\_set}.  Four variants
share a common \texttt{DSU} module interface: a sequential baseline, a
coarse-grained mutex, a per-node lock, and the CAS implementation.
On a workload of 200\,000 nodes and 1\,000\,000 interleaved operations, CAS
peaks at four OCaml domains with a 1.22$\times$ speedup for 20\% unions and a
1.76$\times$ speedup for 50\% unions.  CAS failure rates remain below 0.1\% at
all domain counts, indicating that the performance ceiling is caused by atomic
memory traffic and scheduler overhead rather than by retry storms.
\end{abstract}

% =============================================================================
\section{Goals}
\label{sec:goals}
% =============================================================================

The goal of this project was to implement and evaluate concurrent union-find in
OCaml~5.  Union-find maintains a partition of a fixed element set and supports
\texttt{find} (return a component representative), \texttt{union} (merge two
components), and \texttt{same\_set} (test connectivity).  The structure is a
good fit for a concurrency project because small writes to shared parent
pointers determine global connectivity, so correctness depends on careful
synchronisation rather than just wrapping a sequential algorithm.

\textbf{Research question.} Does a CAS-based concurrent union-find in OCaml~5
outperform lock-based alternatives on realistic interleaved graph-connectivity
workloads, and how does graph structure (density, degree distribution) affect
CAS contention and speedup?
% =============================================================================
\section{Background}
\label{sec:background}
% =============================================================================

Sequential union-find represents each component as a rooted tree in which every
element stores a parent index (roots point to themselves).  Union-by-rank links
the shallower tree under the deeper one, incrementing the winner's rank when
both trees have equal rank.  Path compression shortens future searches by
rerouting nodes directly toward the root.  Together these optimisations give
near-constant amortised cost per operation~\cite{AoMPP}.

In the concurrent setting every optimisation becomes a potential race.
A coarse mutex serialises all operations: correct but unscalable.  Per-node
locks allow concurrent access to disjoint subtrees but require a fixed lock
ordering to avoid deadlock.  CAS uses optimistic synchronisation: read the
current state, compute an update, and install it only if the field still holds
the observed value; retry on failure.

The CAS approach suits union-find because the critical writes are narrow
(one parent index, one rank integer) and failures are rare.  Our implementation
uses \emph{path halving}: when traversing from node $x$, if $x$'s grandparent
$g$ differs from its parent $p$, attempt \texttt{compare\_and\_set parent[x] p g}
and continue from $p$ regardless of whether the CAS succeeded.  For
\texttt{union}, the critical step is:

\begin{lstlisting}[language=caml,caption={CAS link in \texttt{union} (\texttt{cas\_union\_find.ml}).},label={lst:cas-union}]
(* winner/loser chosen by rank, then by smaller index *)
if Atomic.compare_and_set t.parent.(loser) loser winner then (
  (* optional rank increment if equal-rank merge *)
  if same_rank && Atomic.get t.parent.(winner) = winner then begin
    let r = Atomic.get t.rank.(winner) in
    ignore (Atomic.compare_and_set t.rank.(winner) r (r + 1))
  end;
  true)
else
  attempt ()   (* loser is no longer a root; retry *)
\end{lstlisting}

The linearisation point of a successful \texttt{union} is the CAS that sets
\texttt{parent[loser] := winner}.  A failed CAS means another domain has
already merged or compressed that root; the operation simply retries with fresh
roots.  A failed rank increment does not break correctness because ranks are
only a balancing heuristic.

% =============================================================================
\section{Tasks Undertaken}
\label{sec:tasks}
% =============================================================================

\subsection{Implementation}

All four implementations satisfy the \texttt{DSU} module type in
\texttt{union\_find.mli}, which exposes \texttt{create}, \texttt{find},
\texttt{union}, \texttt{same\_set}, \texttt{component\_count}, \texttt{stats},
and \texttt{reset\_stats}.  This made tests and benchmarks implementation-agnostic.

\begin{description}[leftmargin=1.5em,labelindent=0em]
  \item[\texttt{sequential\_union\_find.ml}] \texttt{int array} for parent and
    rank; full path compression (two-pass, extra cache misses on the return
    pass) and union-by-rank; no synchronisation.  Ground-truth reference for
    all cross-checks.

  \item[\texttt{mutex\_union\_find.ml}] Single \texttt{Mutex.t} wrapping the
    whole structure.  Every operation acquires it on entry.  Trivially correct;
    eliminates all parallelism.

  \item[\texttt{node\_lock\_union\_find.ml}] One \texttt{Mutex.t} per node;
    \texttt{Atomic.t} arrays for parent and rank.  \texttt{find} reads
    atomically without locking.  \texttt{union} locks both roots in ascending
    index order (deadlock prevention) and re-checks they are still roots before
    linking.  Key subtlety: same-set must be checked \emph{before} acquiring
    locks --- otherwise two domains on the same component deadlock trying to
    acquire the same lock pair.

  \item[\texttt{cas\_union\_find.ml}] \texttt{int Atomic.t array} only; no
    locks.  Uses \emph{path halving} (single forward pass, skips every other
    node) --- avoids the reverse write-back cache misses of full compression and
    fits naturally into a CAS loop.  Failed CAS $\Rightarrow$ retry from fresh
    roots; no domain ever blocks another~\cite{AlistarhFK19}.
\end{description}

\Cref{tab:impl-compare} summarises the key design choices across the four variants.

\begin{table}[H]
\centering
\small
\begin{tabular}{@{}lllll@{}}
\toprule
 & Sequential & Mutex & Node-lock & CAS \\
\midrule
Storage & \texttt{int array} & \texttt{int array} & \texttt{int Atomic.t array} & \texttt{int Atomic.t array} \\
Synchronisation & none & 1 global lock & 1 lock per node & lock-free CAS \\
Blocking? & n/a & yes & yes & no \\
Path compression & full & full & halving & halving \\
Thread-safe & no & yes & yes & yes \\
\bottomrule
\end{tabular}
\caption{Design comparison of the four union-find implementations.}
\label{tab:impl-compare}
\end{table}

\subsubsection*{Pseudocode for the three concurrent implementations}

The three listings below show the essential structure of each concurrent
\texttt{find} and \texttt{union} in pseudocode.  Low-level details (stats
counters, \texttt{component\_count} updates) are omitted for clarity.

\begin{lstlisting}[caption={Global-mutex \texttt{find} and \texttt{union}.},label={lst:mutex-pseudo}]
(* One Mutex.t for the entire structure *)
find t x:
  lock(t.mutex)
  root <- follow parent[] until parent[root] = root
  (* full path compression: repoint all nodes to root *)
  compress t.parent x root
  unlock(t.mutex)
  return root

union t x y:
  lock(t.mutex)
  rx <- find_nolock t x;  ry <- find_nolock t y
  if rx = ry then (unlock; return false)
  link by rank (loser.parent := winner)
  unlock(t.mutex)
  return true
\end{lstlisting}

\begin{lstlisting}[caption={Per-node-lock \texttt{find} and \texttt{union}.},label={lst:nodelock-pseudo}]
(* One Mutex.t per node; Atomic.t arrays for parent and rank *)
find t x:
  (* read-only atomic traversal, no locks *)
  while parent[x] != x:
    x <- Atomic.get parent[x]
  return x

union t x y:
  loop:
    rx <- find t x;  ry <- find t y
    if rx = ry then return false
    (* lock both roots in ascending index order to prevent deadlock *)
    lo, hi <- min(rx,ry), max(rx,ry)
    lock(node_lock[lo]);  lock(node_lock[hi])
    (* re-check: another domain may have merged these roots already *)
    if parent[lo] != lo || parent[hi] != hi then
      unlock both; goto loop
    link by rank (loser.parent := winner)
    unlock both
    return true
\end{lstlisting}

\begin{lstlisting}[caption={CAS \texttt{find} (path halving) and \texttt{union}.},label={lst:cas-pseudo}]
(* int Atomic.t arrays for parent and rank; no locks *)
find t x:
  (* path halving: one forward pass, skip every other node *)
  while true:
    p  <- Atomic.get parent[x]
    if p = x then return x          (* x is root *)
    gp <- Atomic.get parent[p]
    if gp != p then
      CAS parent[x] p gp            (* try to skip p; ignore failure *)
    x <- p

union t x y:
  loop:
    rx <- find t x;  ry <- find t y
    if rx = ry then return false
    winner, loser <- choose by rank, then by smaller index
    if CAS parent[loser] loser winner then
      (* try to increment winner's rank if equal-rank merge *)
      if same_rank && Atomic.get parent[winner] = winner then
        CAS rank[winner] r (r+1)    (* failure is safe: rank is a hint *)
      return true
    (* CAS failed: loser was already merged by another domain *)
    goto loop
\end{lstlisting}

\subsection{Testing}

Correctness was checked at several levels.  All four test suites are run with
\texttt{make test} (which calls \texttt{dune build} first); individual suites
can be run with the targets below.

\begin{description}[leftmargin=1.5em,labelindent=0em]
  \item[\texttt{make correctness}] ~\\
    \texttt{dune exec ./Union\_find\_correctness\_test.exe} \\
    Deterministic tests (empty/singleton structures, basic unions, repeated
    unions, self-unions, component counts, fixed graph examples, disconnected
    components, stats counters) run on all four implementations, followed by
    randomised cross-checks: 30 trials each of Erdős--Rényi, star, and
    disconnected graphs comparing each implementation's canonical partition
    against the sequential reference, parallel replay over four domains, and
    20 high-contention stress trials.

  \item[\texttt{make stm}] ~\\
    \texttt{dune exec ./QCheck\_stm.exe} \\
    QCheck-STM models the DSU as a sequence of successful \texttt{union} calls
    and checks that \texttt{find}/\texttt{union}/\texttt{same\_set} agree with
    the model.  Runs 300 sequential agreement tests and 150 parallel agreement
    tests (with automatic shrinking on failure) for each of mutex, node-lock,
    and CAS.

  \item[\texttt{make lin}] ~\\
    \texttt{dune exec ./QCheck\_lin.exe} \\
    QCheck-Lin generates parallel two-domain histories and checks that some
    sequential interleaving of the observed responses is valid.  Runs 200
    linearisability tests per concurrent implementation over a small universe.

  \item[\texttt{make cas-stress}] ~\\
    \texttt{dune exec ./cas\_stress\_test.exe} \\
    CAS-specific stress: 30 graph-stress trials (random-plus-star edges over
    four domains, partition checked against sequential replay) and 20
    mixed-operation trials (four workers $\times$ 400 operations each,
    counters verified to be non-zero).
\end{description}

Running \texttt{make bench} executes the two performance benchmarks:

\begin{description}[leftmargin=1.5em,labelindent=0em]
  \item[\texttt{dune exec ./union\_find\_bench.exe}] Scaling benchmark across
    1/2/4/8/16 domains for CAS, mutex, and node-lock; outputs
    \texttt{bench\_trend\_results.csv} and the two trend plots.
  \item[\texttt{dune exec ./union\_find\_graph\_bench.exe}] Graph-structure
    benchmark (sparse-random, dense-random, star) for CAS only; outputs
    \texttt{bench\_graph\_structure\_results.csv} and
    \texttt{bench\_graph\_time.png}.
  \item[\texttt{dune exec ./union\_find\_connectivity\_bench.exe}] Connected-component
    labelling benchmark: partitions edges across domains, each domain calls
    \texttt{union} on its slice, then \texttt{component\_count} is verified
    against a sequential reference.  Five graph profiles (sparse-random,
    dense-random, star, grid, chain); outputs
    \texttt{bench\_connectivity\_results.csv} and
    \texttt{bench\_connectivity\_time.png}.
\end{description}

Testing shaped the implementation in several ways: it forced root re-validation
after optimistic reads, correct handling of a failed CAS as a retry signal
rather than an error, and partition comparisons in benchmarks rather than
timing alone.

% =============================================================================
\section{Evaluation}
\label{sec:evaluation}
% =============================================================================

\subsection{Experimental setup}

Benchmarks used 200\,000 nodes and 1\,000\,000 interleaved operations, with one
warm-up run and three timed runs; the median is reported.  Domains were varied
over \{1, 2, 4, 8, 16\}.  The machine was an Intel Core i5-10300H at 2.50\,GHz
with 4 physical cores and 8 hardware threads; 16-domain results are an
oversubscription experiment.  Two operation mixes were evaluated: 20\% unions /
80\% finds and 50\% unions / 50\% finds.  Operations within each run were
interleaved rather than split into a build phase and a query phase, because
separated phases suppress union/find contention.

\subsection{Results}

\Cref{fig:trend20,fig:trend50} show the runtime curves for the two mixes.
\Cref{tab:summary} summarises the best observed median and speedup from
\texttt{bench\_trend\_results.csv}.

\begin{figure}[H]
  \centering
  \includegraphics[width=0.82\linewidth]{bench_trend_20_time.png}
  \caption{Median runtime --- 20\% union / 80\% find workload.}
  \label{fig:trend20}
\end{figure}

\begin{figure}[H]
  \centering
  \includegraphics[width=0.82\linewidth]{bench_trend_50_time.png}
  \caption{Median runtime --- 50\% union / 50\% find workload.}
  \label{fig:trend50}
\end{figure}

\begin{table}[H]
\centering
\begin{tabular}{@{}llrrrr@{}}
\toprule
Workload & Implementation & 1 domain & Best time & Best domains & Speedup \\
\midrule
20\% union & CAS       & 1.081\,s & 0.889\,s & 4  & 1.22$\times$ \\
20\% union & Mutex     & 1.205\,s & 1.205\,s & 1  & ---          \\
20\% union & Node-lock & 4.413\,s & 1.421\,s & 16 & 3.11$\times$ \\
50\% union & CAS       & 1.612\,s & 0.916\,s & 4  & 1.76$\times$ \\
50\% union & Mutex     & 1.284\,s & 1.284\,s & 1  & ---          \\
50\% union & Node-lock & 5.137\,s & 1.920\,s & 16 & 2.68$\times$ \\
\bottomrule
\end{tabular}
\caption{Best observed median runtime and speedup (speedup relative to own 1-domain time; mutex does not scale so no speedup is listed).}
\label{tab:summary}
\end{table}

\textbf{Mutex.} The coarse mutex is competitive at one domain but degrades
steadily as domains increase, reaching roughly $0.35\times$ the single-domain
throughput at 16 domains for both mixes.  The single lock serialises the entire
operation sequence regardless of contention.

\textbf{CAS.} CAS achieves its best time at four domains for both mixes,
matching the physical core count.  Beyond four domains, runtime rises: the
20\% workload slows from 0.889\,s to 1.527\,s at 16 domains; the 50\% workload
from 0.916\,s to 1.096\,s.  However, CAS failure counters show that failures
are not the cause: at 16 domains the 20\% mix recorded only 43 failures out of
$\sim$373\,000 CAS attempts, and the 50\% mix 359 out of $\sim$388\,000 (both
below 0.1\%).  The performance ceiling is therefore caused by atomic memory
traffic, cache-line movement, OCaml domain scheduling, and hardware
oversubscription rather than by retry storms.

\textbf{Node-lock.} Per-node locks are slow at one domain because every
path-traversal step acquires a mutex, but they scale well in relative terms as
independent graph regions can be processed in parallel.  Despite the strong
relative speedup, the absolute best times (1.421\,s and 1.920\,s) are worse
than the best CAS times, so fine-grained locking pays a high constant overhead
that CAS avoids.

\textbf{Graph structure.} \Cref{fig:graph} isolates the CAS implementation
across three graph profiles.  Sparse-random graphs (low contention) show the
clearest benefit at four domains (1.188\,s $\to$ 0.838\,s).  Dense-random and
star/hot-set profiles are faster at one domain but scale no better beyond four,
because the star workload concentrates unions on a small set of representatives,
causing repeated root contention rather than useful path compression.

\begin{figure}[H]
  \centering
  \includegraphics[width=0.82\linewidth]{bench_graph_time.png}
  \caption{CAS runtime across three graph profiles (sparse random, dense random, star).}
  \label{fig:graph}
\end{figure}

\subsection{Connected-component labelling benchmark}

\texttt{union\_find\_connectivity\_bench.exe} evaluates a pure union-only
workload: build a graph, partition its edges across $d$ domains, let each
domain call \texttt{union} on its slice, then read \texttt{component\_count}.
The final count is verified against a sequential reference run.  This isolates
the cost of \emph{processing edges} from the mixed find/union traffic of the
earlier benchmarks.  Five graph profiles are tested; \cref{fig:connectivity}
shows the results.

\begin{figure}[H]
  \centering
  \includegraphics[width=0.82\linewidth]{bench_connectivity_time.png}
  \caption{CAS connected-component labelling across five graph profiles.}
  \label{fig:connectivity}
\end{figure}

\textbf{Why dense-random is much slower.}  Dense-random uses average degree~40
(vs.\ degree~4 for sparse-random), giving 4\,000\,000 edges against
$\approx$400\,000 for all other profiles.  Every edge still requires two
\texttt{find} calls and a potential CAS, so the work is simply 10$\times$
larger.  Beyond that, once the dense graph collapses into a single component
early in the run, every subsequent edge triggers a redundant connectivity check:
both endpoints find the same root, no CAS is attempted, but the full
path-halving traversal is still paid.  The result is $\sim$4--9\,s versus
$\lesssim$1\,s for all other profiles.

\textbf{Scaling by profile.}
\begin{itemize}[noitemsep]
  \item \textbf{Sparse-random, grid} --- moderate scaling up to 4--8 domains;
    disjoint subtrees allow genuine parallelism.
  \item \textbf{Star} --- near-flat or slightly improving; all unions race to
    set the same root (node~0), so contention is high but path lengths are
    trivially short, keeping absolute time low ($\lesssim$0.2\,s).
  \item \textbf{Chain} --- flat across all domain counts; edges are sequential
    by nature ($0$--$1$, $1$--$2$, \ldots), so partitioning gives each domain
    a contiguous run with almost no cross-partition contention.
  \item \textbf{Dense-random} --- rises sharply from 4 to 8 domains (4.2\,s
    $\to$ 8.8\,s) then levels off; the graph is already fully connected after
    the first few thousand edges, so extra domains spend most of their time on
    redundant finds with no useful union work to parallelise.
\end{itemize}

% =============================================================================
\section{Reflection on the Use of LLMs}
\label{sec:llm}
% =============================================================================

We used two LLMs during this project: \textbf{GPT-4.5} (via Codex, used
primarily inside the editor for code assistance) and \textbf{Claude~Sonnet~4}
(used for design discussions, debugging explanations, and report drafting).

\paragraph{What worked.}
Both models were effective at generating boilerplate: module wrappers, stats
counter scaffolding, CSV summary printers, and QCheck property skeletons.
Claude suggested that interleaved operation arrays give a fairer concurrency
benchmark than separated union and find phases (separated phases suppress
union/find contention), which we adopted.  GPT-4.5 helped untangle recursive
lock-acquisition code and suggested the ascending-index locking discipline for
the node-lock implementation.

\paragraph{What did not work.}
Generated explanations were sometimes overconfident about lock-free progress
guarantees.  Our implementation uses optimistic CAS updates validated by tests,
not a formal proof of lock-freedom; we had to be careful not to inherit that
framing.  The models also predicted that higher domain counts would produce
more CAS retries and proportionally slower runtimes, which the failure counters
contradicted --- failures stayed below 0.1\% even at 16 domains.

\paragraph{A concrete bug caught without LLM help.}
GPT-4.5 helped with code involving locks and recursion, but it missed a subtle
deadlock in an early draft of the node-lock \texttt{union}: the code acquired
locks on both roots without first checking whether they were already in the same
set.  Two domains calling \texttt{union} on nodes from the same component would
each identify the same two roots, then race to acquire the same pair of locks,
and deadlock.  The fix --- check \texttt{find(x) = find(y)} before acquiring
any lock, and re-verify root identity after acquiring --- had to be reasoned
through manually from first principles.

\paragraph{What was difficult.}
The hard reasoning was not mechanical: deciding what each operation is allowed
to conclude from an evolving parent forest, identifying safe retry points, and
understanding why a failed path-halving CAS is harmless (the path still shortens
on the next access) whereas a failed union-link CAS requires a full retry from
fresh roots.  These required careful reading of test failures and invariant
analysis that the LLMs could assist with but not replace.

% =============================================================================
\section{Conclusions}
\label{sec:conclusions}
% =============================================================================

A CAS-based union-find in OCaml~5 is correct under the tested concurrent
histories and outperforms both lock-based alternatives at low to moderate domain
counts.  Its peak is at four domains (matching the physical core count), where
it achieves speedups of 1.22$\times$ and 1.76$\times$ for 20\% and 50\% union
mixes respectively.  Beyond four domains, performance declines despite negligible
CAS failure rates, pointing to atomic memory traffic and domain scheduling as
the bottleneck on this hardware.

The per-node-lock design achieves stronger relative speedup but worse absolute
time; the coarse mutex neither scales nor competes.  Future work includes a
formal linearisability argument, experiments on a larger many-core machine, and
embedding the DSU inside a real graph algorithm such as parallel Kruskal's MST.

% =============================================================================
\section*{Contributions}
% =============================================================================
\begin{center}
\begin{tabular}{@{}lcl@{}}
\toprule
\textbf{Member} & \textbf{Contribution (\%)} & \textbf{Main responsibilities} \\
\midrule
Yaswanth Srinivas & 33\% & Benchmark harness, plots, evaluation, report writing \\
Prince            & 33\% & Core implementation (CAS, mutex, node-lock, sequential) \\
Mukesh            & 33\% & Explanation, analysis, testing, QCheck suites \\
Claude Sonnet 4   & 0.5\% & Report drafting, design discussions, debugging explanations \\
GPT-4.5 (Codex)   & 0.5\% & Boilerplate generation, lock-acquisition code assistance \\
\midrule
\textbf{Total} & \textbf{100\%} & \\
\bottomrule
\end{tabular}
\end{center}
 

% =============================================================================
\bibliographystyle{plainurl}
\bibliography{references}
% =============================================================================

\end{document}